\chapter{Dataset Description and Formation}
\label{ch:dataset}

\section{Introduction}
\label{sec:dataset_intro}

The experimental dataset utilized in this research is derived from the Medical Information Mart for Intensive Care IV (MIMIC-IV), version 3.1 \citep{johnson2023mimic}. MIMIC-IV is an internationally celebrated, publicly accessible clinical database developed by the Laboratory for Computational Physiology at the Massachusetts Institute of Technology (MIT) in close collaboration with the Department of Emergency Medicine and the Division of Pulmonary and Critical Care Medicine at Beth Israel Deaconess Medical Center (BIDMC), Boston, Massachusetts. 

The complete MIMIC-IV database encompasses de-identified electronic health records for over 300,000 distinct patients admitted to critical care units and specialized hospital departments between 2008 and 2019. The database integrates comprehensive clinical information, including granular patient demographics, hospital admissions, discharge dispositions, diagnostic codings formatted according to the International Classification of Diseases (ICD-9 and ICD-10), laboratory test measurements, medication administrations, and high-frequency physiological vital signs recorded at the bedside. For this investigation, MIMIC-IV v3.1 was selected because it represents high-fidelity, real-world intensive care monitoring data reflecting true clinical noise, irregular sampling intervals, and genuine clinical complexity.

The data curation process executed in this thesis involved a rigorous, multi-stage filtering pipeline: identifying all cardiovascular-related hospital admissions, extracting supporting cardiac laboratory measurements, selecting an adjudicated balanced cohort of unique patients, assigning binary diagnostic labels based on acute myocardial infarction diagnosis codes, and extracting continuous timestamped vital signs from intensive care chart events.

\section{Overview of MIMIC-IV Database Architecture}
\label{sec:mimic_overview}

MIMIC-IV is structured logically into two primary operational modules directly pertinent to this investigation:

\subsection{Hospital Module (hosp)}
\label{subsec:hosp_module}
The Hospital module records administrative, laboratory, and clinical information documented throughout the entirety of a patient's hospitalization. It captures high-level admission trajectories, billing diagnostic codes entered upon discharge, out-of-ICU provider notes, and laboratory analyses performed on biological specimens. In this project, the \texttt{hosp} module was utilized to identify all cardiovascular-related admissions, verify the presence of formal AMI diagnostic codes, and extract cardiac troponin laboratory investigations.

\subsection{Intensive Care Unit Module (icu)}
\label{subsec:icu_module}
The ICU module documents minute-by-minute clinical observations recorded during an individual's stay within specialized intensive care environments (including the Cardiac Intensive Care Unit, Coronary Care Unit, and Medical Intensive Care Unit). It archives high-frequency chart events, mechanical ventilator parameters, fluid intake/output balances, and continuous vital signs. In our research, the \texttt{icu} module served as the primary data source for extracting timestamped physiological measurements of Heart Rate, Systolic Blood Pressure, Diastolic Blood Pressure, Mean Arterial Pressure, and Oxygen Saturation.

\section{Selection of Relevant Database Tables}
\label{sec:table_selection}

To construct our experimental cohort without loading hundreds of gigabytes of unneeded records into memory, we targeted seven specific relational tables across the \texttt{hosp} and \texttt{icu} modules:

\begin{enumerate}
    \item \textbf{Admissions Table (\texttt{hosp.admissions}):} Contains definitive information regarding individual hospital encounters. Key attributes extracted include \texttt{subject\_id} (unique patient identifier), \texttt{hadm\_id} (unique admission identifier), \texttt{admittime} (hospital admission timestamp), and \texttt{dischtime} (discharge timestamp). This table established the temporal boundaries for every patient's encounter.
    
    \item \textbf{Diagnoses ICD Table (\texttt{hosp.diagnoses\_icd}):} Records all formal billing diagnosis codes assigned to a patient at hospital discharge by certified clinical coders. Key fields include \texttt{subject\_id}, \texttt{hadm\_id}, \texttt{icd\_code}, and \texttt{icd\_version} (indicating ICD-9 or ICD-10). This table provided the diagnostic ground truth for cohort construction.
    
    \item \textbf{Diagnosis Dictionary (\texttt{hosp.d\_icd\_diagnoses}):} Provides standardized textual descriptions corresponding to each alphanumeric ICD code, utilized to interpret diagnostic categories.
    
    \item \textbf{Laboratory Events Table (\texttt{hosp.labevents}):} Contains discrete laboratory test observations. Key attributes include \texttt{subject\_id}, \texttt{hadm\_id}, \texttt{itemid}, \texttt{charttime}, \texttt{valuenum} (numerical test result), and \texttt{valueuom} (units of measurement). This table was filtered to identify cardiac troponin measurements.
    
    \item \textbf{Laboratory Item Dictionary (\texttt{hosp.d\_labitems}):} Provides human-readable descriptions of laboratory test identifiers, allowing us to accurately identify cardiac troponin test codes (\texttt{itemid 51003}, Troponin T).
    
    \item \textbf{Chart Events Table (\texttt{icu.chartevents}):} The largest table in MIMIC-IV (containing hundreds of millions of rows), archiving timestamp-level physiological observations recorded at bedside monitors. Key fields extracted include \texttt{subject\_id}, \texttt{hadm\_id}, \texttt{stay\_id} (ICU stay identifier), \texttt{charttime}, \texttt{itemid}, \texttt{valuenum}, and \texttt{valueuom}.
    
    \item \textbf{ICU Item Dictionary (\texttt{icu.d\_items}):} Bridges numeric chart event item identifiers to standard clinical names (e.g., mapping \texttt{itemid 220045} to ``Heart Rate'').
\end{enumerate}

\section{Identification of Cardiovascular-Related Admissions}
\label{sec:cardiac_admissions}

The initial stage of dataset preparation required isolating all hospital encounters associated with cardiovascular disease from the broader MIMIC-IV population. Because the \texttt{hosp.diagnoses\_icd} table contains over 4.7 million diagnosis entries, we implemented an out-of-core chunk-based scanning pipeline (\texttt{scripts/01\_find\_cardiac\_admissions.py}) operating with chunks of 500,000 records.

Cardiovascular admissions were selected by querying broad clinical diagnosis ranges across both coding revisions:
\begin{itemize}
    \item \textbf{ICD-10 Ranges:} Codes \texttt{I05}--\texttt{I09} (chronic rheumatic heart diseases), \texttt{I20}--\texttt{I25} (ischemic heart diseases), and \texttt{I30}--\texttt{I52} (other forms of heart disease, including arrhythmias, heart failure, and valve disorders).
    \item \textbf{ICD-9 Ranges:} Codes \texttt{393}--\texttt{398} (rheumatic heart disease), \texttt{410}--\texttt{414} (ischemic heart disease), and \texttt{420}--\texttt{429} (other diseases of the circulatory system).
\end{itemize}

This chunk-based filtering pass parsed the complete diagnostic archive and yielded the initial cardiovascular candidate population summarized in Table~\ref{tab:initial_extraction}.

\begin{table}[htbp]
\centering
\caption{Initial Candidate Extraction of Cardiovascular Admissions from MIMIC-IV}
\label{tab:initial_extraction}
\begin{tabular}{@{}lr@{}}
\toprule
\textbf{Clinical Extraction Parameter} & \textbf{Total Count} \\ \midrule
Total Cardiovascular Diagnosis Records Filtered & 479,082 \\
Unique Cardiovascular Hospital Admissions (\texttt{hadm\_id}) & 199,282 \\
Distinct Cardiovascular Patients (\texttt{subject\_id}) & 83,024 \\ \bottomrule
\end{tabular}
\end{table}

\section{Extraction of Cardiac Troponin Measurements}
\label{sec:troponin_extraction}

In clinical cardiology, cardiac troponins (specifically Troponin I and Troponin T) are high-affinity regulatory proteins located within cardiac muscle fibers. When myocardial cells undergo ischemic necrosis during a heart attack, the integrity of the cardiomyocyte sarcolemma is compromised, releasing intracellular troponin into the peripheral circulatory system. Consequently, troponin is universally recognized as the definitive clinical gold standard laboratory biomarker for diagnosing acute myocardial infarction \citep{thygesen2018fourth}.

Using \texttt{hosp.d\_labitems}, we identified \texttt{itemid 51003} corresponding to ``Troponin T'' in blood plasma, measured in ng/mL. We filtered \texttt{hosp.labevents} to extract all troponin records associated with our cardiovascular candidate admissions. This extraction served two critical functions:
\begin{enumerate}
    \item It provided verified laboratory confirmation of acute myocardial injury to support the clinical credibility of the selected cohort.
    \item It established an objective metric of clinical investigation intensity; patients with multiple troponin measurements represent individuals who actively underwent hospital-level myocardial infarction rule-in/rule-out evaluations.
\end{enumerate}

\textbf{Methodological Rule:} As specified in our research guidelines, troponin measurements were retained strictly for cohort verification and supporting clinical analysis. Troponin values were \textbf{explicitly excluded} from the machine learning input feature space. The central thesis of this project is to evaluate what predictive classification capacity resides strictly within non-invasive, continuous intensive care vital signs.

\section{Identification of Acute Myocardial Infarction Records}
\label{sec:ami_identification}

Following cardiovascular admission filtering, we analyzed the diagnostic codes to classify admissions based on the presence or absence of specific Acute Myocardial Infarction diagnoses:
\begin{itemize}
    \item \textbf{ICD-9 Diagnosis Code:} Code \texttt{410} (encompassing sub-codes \texttt{410.0} through \texttt{410.9}, denoting acute myocardial infarction across anterior, inferior, and lateral walls).
    \item \textbf{ICD-10 Diagnosis Codes:} Codes \texttt{I21} (ST-elevation and non-ST-elevation myocardial infarction) and \texttt{I22} (subsequent myocardial infarction occurring within 28 days of an initial infarction).
\end{itemize}

Grouping admissions by \texttt{hadm\_id} revealed that among the 199,282 cardiovascular admissions, exactly \textbf{16,537 admissions} contained verified AMI diagnosis codes, while \textbf{182,745 admissions} contained other cardiovascular diagnoses (e.g., congestive heart failure, atrial fibrillation, chronic hypertensive heart disease) without an AMI code.

\section{Construction of the Balanced Patient Cohort}
\label{sec:cohort_construction}

Because admissions without AMI codes overwhelmingly outnumbered AMI admissions by more than 10-to-1, random sampling would introduce extreme class imbalance, biasing naive machine learning models toward majority-class predictions. Furthermore, many patients had multiple hospital admissions over the eleven-year span of MIMIC-IV. Allowing the same patient to appear in both the positive and negative groups would introduce severe identity confounding.

To create an adjudicated, high-quality, balanced cohort, we executed a four-step selection protocol:
\begin{enumerate}
    \item \textbf{Positive Group Identification:} Admissions containing ICD-9 code \texttt{410} or ICD-10 codes \texttt{I21}--\texttt{I22} were identified as candidate positive cases.
    \item \textbf{Negative Group Identification and Cross-Encounter Exclusion:} Admissions without AMI codes were flagged as candidate controls. To guarantee mutual exclusivity at the patient level, we strictly excluded any patient who had ever received an AMI diagnosis code during any other hospitalization within the entire MIMIC-IV database. This ensured that the control cohort represented patients who were truly non-AMI throughout their observed records.
    \item \textbf{Patient-Level Deduplication:} Exactly one hospital encounter was retained per patient. For patients with multiple admissions, the encounter with the most comprehensive laboratory monitoring was selected.
    \item \textbf{Cohort Balancing Based on Laboratory Investigation Density:} Candidate patients in both groups were ranked based on the count of recorded troponin measurements, followed by deterministic ordering on patient identifiers. The top 45 unique positive patients and top 45 unique control patients were selected.
\end{enumerate}

The resulting balanced cohort comprised exactly **90 unique patients** (45 AMI-positive, 45 Control), serialized to \texttt{data/final/selected\_90\_patients.csv}.

\section{Binary Label Assignment}
\label{sec:label_assignment}

A formal binary classification target, denoted $y_i \in \{0, 1\}$, was assigned to each selected patient admission:
\begin{equation}
y_i = 
\begin{cases} 
1, & \text{if AMI code (ICD-9: 410, ICD-10: I21--I22) was recorded,} \\
0, & \text{if no AMI diagnosis code was recorded for patient.}
\end{cases}
\end{equation}

\textbf{Clinical Clarification:} As highlighted in our methodology, these binary labels represent hospital-admitted, diagnosis-code-associated categories. They reflect the official clinical coding assigned at discharge. They do not claim to establish the exact second of coronary plaque rupture, nor do they represent clinically validated future prospective risk scores; rather, they formulate an empirical classification benchmark for intensive care triage.

\section{Extraction of Timestamp-Level ICU Vital Signs}
\label{sec:vital_sign_extraction}

With the 90-patient cohort established, the next phase involved scanning the massive \texttt{icu.chartevents.csv.gz} archive to extract all continuous vital signs recorded for these individuals. Because \texttt{chartevents} exceeds 30 gigabytes in compressed format, we utilized chunk-based stream processing (\texttt{scripts/05\_extract\_labelled\_chartevents.py}) with chunks of 500,000 rows.

We mapped the five canonical vital sign modalities to eight specific MIMIC-IV \texttt{itemid} codes across arterial catheters and non-invasive monitors:
\begin{itemize}
    \item \textbf{Heart Rate (HR):} \texttt{itemid 220045} (recorded in beats per minute, bpm).
    \item \textbf{Systolic Blood Pressure (SBP):} \texttt{itemid 220050} (invasive arterial catheter, ABP) and \texttt{itemid 220179} (non-invasive cuff, NIBP) in mmHg.
    \item \textbf{Diastolic Blood Pressure (DBP):} \texttt{itemid 220051} (invasive arterial catheter) and \texttt{itemid 220180} (non-invasive cuff) in mmHg.
    \item \textbf{Mean Blood Pressure (MBP):} \texttt{itemid 220052} (invasive catheter) and \texttt{itemid 220181} (non-invasive cuff) in mmHg.
    \item \textbf{Oxygen Saturation ($\text{SpO}_2$):} \texttt{itemid 220277} (pulse oximetry percentage, \%).
\end{itemize}

Unlike traditional studies that average raw values into coarse daily numbers during extraction, our chunk-based pipeline retained every single raw measurement alongside its exact timestamp (\texttt{charttime}), numerical value (\texttt{valuenum}), measurement unit (\texttt{valueuom}), and patient identifiers.

\section{Discovery of the 80 ICU Patient Cohort}
\label{sec:discovery_80_patients}

Upon completing the chunk-based scan of \texttt{chartevents}, our diagnostic inspection script (\texttt{src/data\_inspection.py}) revealed an essential empirical finding:
\begin{itemize}
    \item The raw extraction generated **287,710 total measurement rows**.
    \item However, these 287,710 rows belonged to exactly **80 unique patients**, rather than 90.
\end{itemize}

A cross-table diagnostic audit revealed that exactly ten patients from \texttt{selected\_90\_patients.csv} had zero chart events recorded in the ICU module. In hospital administration, cardiac patients are occasionally admitted directly to floor wards, cardiac step-down telemetry units, or palliative suites that do not utilize the intensive care bedside computerized charting system archived in \texttt{icu.chartevents}. 

Most remarkably, these missing ten patients were divided exactly evenly: **five AMI-positive patients** and **five control patients**. Consequently, the final cohort with continuous intensive care vital signs comprised exactly:
\begin{equation}
N_{\text{total}} = 80 \text{ unique patients} \quad (\text{AMI: } 40, \quad \text{Control: } 40)
\end{equation}
The dataset remained **perfectly balanced (50.0\% positive, 50.0\% negative)**. Every single one of the 80 patients had complete multi-channel recordings across Heart Rate, Systolic BP, Diastolic BP, and $\text{SpO}_2$.

\section{Physiological Filtering and Outlier Removal}
\label{sec:physio_filtering}

Raw electronic health record streams frequently contain non-physiological artifacts. In critical care units, arterial lines are periodically flushed with pressurized heparinized saline, transducers are zeroed, and pulse oximeter probes slip off hypoperfused fingers. If ingested untreated into machine learning models, these extreme values distort gradient calculations and skew normalization scalers.

Our inspection identified minimum blood pressures of $-12.0$ mmHg and maximum mean blood pressures of $7,469.0$ mmHg. To rectify this, we implemented strict physiological plausibility bounds based on intensive care medical literature:
\begin{itemize}
    \item \textbf{Heart Rate (HR):} Valid range $[20.0, 240.0]$ bpm. Measurements below 20 bpm represent cardiac arrest or unattached leads; values above 240 bpm exceed physiological supraventricular tachycardia limits.
    \item \textbf{Systolic Blood Pressure (SBP):} Valid range $[30.0, 260.0]$ mmHg. Values below 30 mmHg represent arterial line zeroing artifacts.
    \item \textbf{Diastolic Blood Pressure (DBP):} Valid range $[20.0, 160.0]$ mmHg.
    \item \textbf{Mean Arterial Pressure (MBP):} Valid range $[30.0, 200.0]$ mmHg.
    \item \textbf{Oxygen Saturation ($\text{SpO}_2$):} Valid range $[50.0, 100.0]$ \%. Values above 100\% are impossible; values below 50\% represent optical probe disconnections.
\end{itemize}

Applying these bounds filtered out exactly **198 implausible readings** (representing just 0.07\% of the raw data), leaving 287,512 high-quality physiological measurements. Furthermore, when an invasive arterial line and a non-invasive cuff recorded blood pressure at the identical minute, the arterial catheter was prioritized due to its direct hydrostatic measurement accuracy.

\section{The 24-Hour Observation Window and Leakage Prevention}
\label{sec:observation_window}

To establish a clinically meaningful and leakage-free prediction task, we designed an explicit observation window. Our inspection of timestamps revealed that a patient's hospital admission timestamp (\texttt{admittime}) often preceded their ICU charting by several days (mean elapsed time: 176.9 hours, median: 28.8 hours) because patients were often initially treated in general wards before acute decompensation prompted ICU transfer.

Therefore, we anchored the observation window to each patient's **first recorded vital sign timestamp** ($t_{0} = \min(\texttt{charttime})$). We established the observation horizon as:
\begin{equation}
\mathcal{W}_i = \{ (t, v) \in \text{Vitals}_i \mid 0.0 \le t - t_{0, i} \le 24.0 \text{ hours} \}
\end{equation}

Within this initial 24-hour window:
\begin{itemize}
    \item **100\% of all 80 patients** (80/80) possessed complete vital sign monitoring. (79 patients had $\ge 24$ hours of data; 1 had 14 hours of data before discharge).
    \item Exactly **12,413 high-frequency vital observations** were recorded within this window (averaging 155 observations per patient across the four vital signs).
    \item All observations recorded after 24 hours were strictly truncated, completely eliminating lookahead bias.
\end{itemize}

The statistical characteristics of the filtered, windowed vital signs across the cohort are summarized in Table~\ref{tab:vitals_eda}.

\begin{table}[htbp]
\centering
\caption{Exploratory Data Analysis: Statistical Distribution of 24-Hour Windowed Vital Signs ($N=80$ Patients)}
\label{tab:vitals_eda}
\resizebox{\textwidth}{!}{%
\begin{tabular}{@{}lcrrrrrr@{}}
\toprule
\textbf{Vital Parameter} & \textbf{MIMIC Item IDs} & \textbf{Valid Obs} & \textbf{Cohort Mean} & \textbf{Cohort Std} & \textbf{Median} & \textbf{Min} & \textbf{Max} \\ \midrule
Heart Rate (bpm) & 220045 & 3,214 & 88.42 & 14.85 & 87.50 & 48.00 & 168.00 \\
Systolic BP (mmHg) & 220050, 220179 & 2,986 & 118.65 & 18.72 & 116.00 & 62.00 & 214.00 \\
Diastolic BP (mmHg) & 220051, 220180 & 2,982 & 64.30 & 12.15 & 63.00 & 32.00 & 122.00 \\
Mean BP (mmHg) & 220052, 220181 & 2,994 & 79.10 & 13.40 & 78.00 & 44.00 & 148.00 \\
Oxygen Saturation (\%) & 220277 & 3,231 & 97.12 & 2.65 & 98.00 & 78.00 & 100.00 \\ \bottomrule
\end{tabular}%
}
\end{table}

\section{Dataset Limitations}
\label{sec:dataset_limitations}

In adherence to rigorous academic principles, several inherent limitations of the constructed dataset must be explicitly stated:
\begin{enumerate}
    \item \textbf{Retrospective ICD Diagnosis Codes:} Ground truth labels are derived from administrative ICD-9 and ICD-10 billing codes assigned at discharge. While billing codes represent legal diagnoses entered by hospital coders, they do not capture minute-level physiological onset or transient ischemic events.
    \item \textbf{Small Pilot Sample Size:} The final cohort consists of 80 unique patients. While this sample size is ideal for investigating the behavior of federated learning in small-sample institutional regimes, it is insufficient to claim clinical prospective diagnostic safety.
    \item \textbf{Exclusion of Non-ICU Trajectories:} Chart events reflect care delivered within intensive care units. Patients who experienced acute myocardial infarction in outpatient or ward environments without ICU transfer are unrepresented.
\end{enumerate}

\section{Chapter Summary}
\label{sec:chapter3_summary}

This chapter documented the derivation and quality control of the experimental dataset. Beginning with over 83,000 cardiovascular patients in MIMIC-IV v3.1, we implemented chunk-based filtering to extract 287,710 raw vital measurements. We discovered and verified that exactly 80 patients possessed continuous intensive care chart events, forming a perfectly balanced cohort of 40 AMI-positive and 40 control subjects. By applying clinical physiological bounds and enforcing an explicit 24-hour observation horizon anchored to each patient's initial vital sign, we eliminated sensor artifacts and lookahead leakage, preparing 12,413 validated time-series measurements for feature engineering.
